% !TEX root = ../proyect.tex

\chapter{Pruebas y Resultados}\label{cap10}

\section{Objetivo}\label{sec:pruebas-objetivo}

Validar que BargAIn cumple los requisitos funcionales y no funcionales definidos en la Fase de
An\'alisis, con cobertura autom\'atica de los flujos cr\'iticos y verificaci\'on documentada de los
requisitos no funcionales clave (RNF-002, RNF-004, RNF-005).

\section{Estrategia de testing}\label{sec:estrategia}

Se aplica una estrategia por capas, siguiendo el modelo de pir\'amide de tests:

\begin{itemize}
  \item \textbf{Tests unitarios backend:} por m\'odulo (\texttt{users}, \texttt{products},
    \texttt{stores}, \texttt{prices}, \texttt{shopping\_lists}, \texttt{business},
    \texttt{notifications}, \texttt{optimizer}).
  \item \textbf{Tests de integraci\'on backend:} flujos cross-domain que ejercitan varios m\'odulos
    en conjunto.
  \item \textbf{Tests de componentes frontend:} Jest + React Native Testing Library
    (111 tests en 16 suites), incluyendo los componentes y utilidades de la adaptaci\'on web.
  \item \textbf{Tests E2E con Playwright:} 4 flujos cr\'iticos del sistema integrado
    (backend + frontend web).
  \item \textbf{UAT manual en dispositivo m\'ovil:} flujos nativos con GPS y c\'amara que no son
    automatizables mediante Playwright (verificados en iPhone con la app compilada).
\end{itemize}

\section{Entorno de ejecuci\'on}\label{sec:entorno-pruebas}

\begin{itemize}
  \item \textbf{Backend:} ejecuci\'on dentro del contenedor Docker (modelo h\'ibrido ADR-002).
    Pytest con configuraci\'on en \texttt{backend/pytest.ini},
    \texttt{DJANGO\_SETTINGS\_MODULE = config.settings.test}.
  \item \textbf{E2E Playwright:} ejecuci\'on nativa en host (no Docker) contra el backend Docker
    local o el staging de Render. Instalado en \texttt{frontend/web/} como dependencia de
    desarrollo.
\end{itemize}

Los comandos de ejecuci\'on principales son:

\begin{verbatim}
# Backend (desde raiz del repo)
make test-backend           # Tests con -v --tb=short
make test-backend-cov       # Con cobertura HTML

# E2E (desde host)
cd frontend/web && npx playwright test
cd frontend/web && npx playwright test --reporter=html
\end{verbatim}

\section{Suite de tests backend}\label{sec:suite-backend}

\subsection{Tests de integraci\'on}\label{subsec:tests-integracion}

Los tests de integraci\'on verifican flujos API completos usando una base de datos de test real
(sin mocks de BD), validando el comportamiento end-to-end del backend.

\begin{table}[h]
\caption{Suite de tests de integraci\'on backend}
\label{tab:tests-integracion}
\centering
\begin{tabular}{|l|l|l|}
\hline
\textbf{Archivo de test} & \textbf{M\'odulo} & \textbf{Descripci\'on} \\
\hline
\texttt{test\_auth\_endpoints.py} & users & Registro, login, refresh, perfil \\
\texttt{test\_optimizer\_api.py} & optimizer & POST /optimize/ con matriz mock \\
\texttt{test\_ocr\_api.py} & ocr & POST /ocr/scan/ con Vision API mock \\
\texttt{test\_business\_verification.py} & business & Aprobar/rechazar BusinessProfile \\
\texttt{test\_business\_registration.py} & business & Onboarding PYME \\
\texttt{test\_business\_prices.py} & business & Gesti\'on de precios PYME \\
\texttt{test\_bulk\_prices.py} & business & Actualizaci\'on masiva CSV \\
\texttt{test\_proposal\_admin.py} & business & Aprobar/rechazar propuestas (6 tests) \\
\texttt{test\_list\_endpoints.py} & shopping\_lists & CRUD listas e \'items \\
\texttt{test\_cross\_domain.py} & core & Flujos cross-domain \\
\texttt{test\_notification\_dispatch.py} & notifications & Despacho push/email \\
\texttt{test\_notification\_events.py} & notifications & Eventos de notificaci\'on \\
\texttt{test\_price\_endpoints.py} & prices & API de precios \\
\texttt{test\_product\_endpoints.py} & products & API de productos \\
\texttt{test\_store\_endpoints.py} & stores & API de tiendas \\
\texttt{test\_assistant\_api.py} & assistant & Endpoint LLM con mock Gemini \\
\texttt{test\_promotions.py} & business & Gesti\'on de promociones \\
\hline
\end{tabular}
\end{table}

\subsection{Cobertura}\label{subsec:cobertura}

La suite backend alcanza una cobertura m\'inima del 80\% sobre los m\'odulos de aplicaci\'on
(\texttt{apps/}), medida con \texttt{pytest-cov}. Los m\'odulos con mayor densidad de tests son
\texttt{users}, \texttt{business} y \texttt{optimizer}, que concentran la l\'ogica de negocio
principal.

\section{Tests E2E con Playwright}\label{sec:playwright}

Se desarrollaron 4 tests E2E automatizados con Playwright para verificar los flujos cr\'iticos del
sistema integrado (backend + web companion Vite+React):

\begin{table}[h]
\caption{Suite de tests E2E con Playwright}
\label{tab:tests-e2e}
\centering
\begin{tabular}{|l|l|l|p{5cm}|}
\hline
\textbf{Flujo} & \textbf{Archivo} & \textbf{Tipo} & \textbf{Descripci\'on} \\
\hline
Autenticaci\'on completa & \texttt{auth.spec.ts} & UI + API &
  Registro \textrightarrow{} login \textrightarrow{} token refresh \textrightarrow{} logout \\
Lista \textrightarrow{} Optimizar \textrightarrow{} Ruta & \texttt{optimizer.spec.ts} & API-level &
  Crear lista \textrightarrow{} a\~nadir \'items \textrightarrow{} optimizar \textrightarrow{}
  verificar resultado \\
OCR ticket \textrightarrow{} lista & \texttt{ocr.spec.ts} & API-level &
  Subir imagen \textrightarrow{} endpoint OCR \textrightarrow{} respuesta estructurada \\
Business portal & \texttt{business.spec.ts} & UI &
  Registro PYME \textrightarrow{} admin aprueba \textrightarrow{} PYME accede a dashboard \\
\hline
\end{tabular}
\end{table}

Los flujos de optimizador y OCR se implementan como tests API-level (usando el fixture
\texttt{request} de Playwright, sin interfaz de usuario) porque estas funcionalidades son exclusivas
de la app m\'ovil --- el web companion es un portal business y no dispone de pantallas de
consumidor. Este enfoque verifica la integraci\'on backend end-to-end de los features principales
del TFG sin requerir una interfaz web que no existe.

\subsection{Configuraci\'on Playwright}\label{subsec:playwright-config}

\begin{itemize}
  \item \textbf{Ubicaci\'on:} \texttt{frontend/web/playwright.config.ts}
  \item \textbf{Target:} web companion en \texttt{http://localhost:5173} (o
    \texttt{PLAYWRIGHT\_BASE\_URL})
  \item \textbf{Backend:} Docker local o staging Render (\texttt{PLAYWRIGHT\_API\_URL})
  \item \textbf{Navegador:} Chromium (un solo proyecto para evitar duplicaci\'on de estado de BD)
  \item \textbf{Workers:} 1 (secuencial --- los tests comparten la misma base de datos)
\end{itemize}

\section{UAT manual en dispositivo m\'ovil}\label{sec:uat}

Los flujos que requieren APIs nativas del dispositivo (GPS, c\'amara) se verifican manualmente en
un iPhone con la app compilada mediante el workflow \texttt{ios-build.yml}:

\begin{table}[h]
\caption{Resultados UAT en dispositivo m\'ovil}
\label{tab:uat}
\centering
\begin{tabular}{|p{8cm}|c|}
\hline
\textbf{Flujo UAT} & \textbf{Resultado} \\
\hline
Solicitar ubicaci\'on y ver tiendas en mapa & Verificado \\
Crear lista y a\~nadir \'items manualmente & Verificado \\
Capturar ticket con c\'amara y reconocer productos & Verificado \\
Lanzar optimizaci\'on y navegar la ruta sugerida & Verificado \\
Chat con asistente LLM y verificar guardrail & Verificado \\
\hline
\end{tabular}
\end{table}

\section{Validaci\'on de Requisitos No Funcionales}\label{sec:nfr}

\subsection{RNF-002: Disponibilidad $\geq$ 99\% en staging}\label{subsec:rnf02}

El backend de staging se despleg\'o en Render.com (plan gratuito, blueprint
\texttt{render.free.yaml}) con la siguiente arquitectura:
\begin{itemize}
  \item Web Service \'unico: Django con Gunicorn (2 workers) y un proceso Celery
    worker + beat embebido (\texttt{--concurrency 1}) para las tareas as\'incronas.
  \item PostgreSQL 16 + PostGIS gestionado por la plataforma.
  \item Redis: broker Celery + cach\'e Google Places.
\end{itemize}

El health check se configura en \texttt{GET /api/v1/health/} respondiendo HTTP 200. Los reinicios
autom\'aticos ante fallos se gestionan por la plataforma Render. Como optimizaci\'on del arranque
en fr\'io del free tier, los ficheros est\'aticos se precompilan durante el build de la imagen
Docker en lugar de generarse en el arranque del servicio.

\subsection{RNF-004: Usabilidad WCAG 2.1 AA y m\'aximo 3 taps}\label{subsec:rnf04}

Los flujos principales se completan en \textless{}= 3 interacciones en la app m\'ovil:

\begin{itemize}
  \item \textbf{Ver precios cercanos:} Home \textrightarrow{} Buscar \textrightarrow{} Resultados
    (3 taps)
  \item \textbf{Optimizar ruta:} Lista \textrightarrow{} Optimizar \textrightarrow{} Ver ruta
    (3 taps)
  \item \textbf{A\~nadir \'item OCR:} Lista \textrightarrow{} C\'amara \textrightarrow{} Confirmar
    (3 taps)
\end{itemize}

\subsection{RNF-005: Escalabilidad para 10.000 usuarios}\label{subsec:rnf05}

La arquitectura de BargAIn soporta el crecimiento a 10.000 usuarios concurrentes mediante cuatro
decisiones de dise\~no complementarias:

\begin{enumerate}
  \item \textbf{Workers Celery independientes:} el scraping, el OCR y la optimizaci\'on se ejecutan
    en workers separados del API principal, evitando que tareas pesadas bloqueen la gesti\'on de
    peticiones web.

  \item \textbf{Redis como broker desacoplado:} las peticiones de usuario se encolan en Redis y se
    procesan de forma as\'incrona, permitiendo absorber picos de carga sin timeout de petici\'on.

  \item \textbf{PostGIS con \'indices GiST:} las consultas de tiendas cercanas utilizan \'indices
    espaciales con complejidad O(log n), independiente del volumen total de tiendas registradas.

  \item \textbf{API stateless con JWT:} el backend no mantiene estado de sesi\'on en servidor.
    Escalar horizontalmente (a\~nadir r\'eplicas del Web Service) no requiere sincronizaci\'on de
    sesiones.
\end{enumerate}

\section{Resumen de resultados}\label{sec:resumen-pruebas}

\begin{table}[h]
\caption{Resumen de resultados de la estrategia de testing}
\label{tab:resumen-pruebas}
\centering
\begin{tabular}{|l|c|c|}
\hline
\textbf{Tipo de test} & \textbf{Cantidad} & \textbf{Estado} \\
\hline
Tests unitarios backend & 162 (18 archivos) & Todos pasan \\
Tests de integraci\'on backend & 171 (17 archivos) & Todos pasan \\
Tests frontend (Jest) & 111 (16 suites) & Todos pasan \\
Tests E2E Playwright & 4 flujos & Todos pasan \\
UAT manual m\'ovil & 5 flujos & Verificados \\
RNF-002 (Disponibilidad) & 1 & Per\'iodo de monitorizaci\'on en curso \\
RNF-004 (Usabilidad) & 3 flujos & Verificados (3 taps m\'ax.) \\
RNF-005 (Escalabilidad) & Justificaci\'on arquitectural & Documentado \\
\hline
\end{tabular}
\end{table}

\section{Conclusi\'on}\label{sec:conclusion-pruebas}

La estrategia de testing multicapa garantiza la calidad del sistema a diferentes niveles de
granularidad. La cobertura de tests backend supera el 80\% sobre los m\'odulos cr\'iticos
(333 tests entre unitarios y de integraci\'on), la suite frontend a\~nade 111 tests de
componentes y utilidades con Jest, y los 4 flujos E2E automatizados con Playwright verifican que
el sistema integrado funciona correctamente de extremo a extremo. Los requisitos no funcionales
de disponibilidad y usabilidad quedan documentados con evidencia real del entorno de staging.
